\documentclass[10pt,a4paper]{article}
\usepackage[T1]{fontenc}
\usepackage[utf8]{inputenc}
\usepackage{mathptmx}
\usepackage[left=2cm,right=2cm,top=1.8cm,bottom=1.8cm]{geometry}
\usepackage{amsmath,amssymb}
\usepackage{graphicx}
\usepackage{booktabs}
\usepackage{float}
\usepackage{array}
\usepackage{microtype}
\usepackage[font=small,labelfont=bf,skip=4pt]{caption}
\usepackage{titlesec}
\usepackage{enumitem}
\usepackage[hidelinks]{hyperref}

\titlespacing*{\section}{0pt}{6pt plus 2pt minus 1pt}{3pt}
\titlespacing*{\subsection}{0pt}{5pt plus 1pt minus 1pt}{2pt}
\titleformat{\section}{\normalfont\large\bfseries}{\thesection}{0.6em}{}
\titleformat{\subsection}{\normalfont\normalsize\bfseries}{\thesubsection}{0.6em}{}
\setlist[itemize]{topsep=2pt,itemsep=1pt,parsep=0pt,leftmargin=1.2em}
\setlength{\abovedisplayskip}{4pt}
\setlength{\belowdisplayskip}{4pt}
\setlength{\textfloatsep}{8pt}
\renewcommand{\arraystretch}{1.05}
\linespread{0.97}\selectfont

\title{\vspace{-1.2cm}\bfseries Tracing the Markowitz Efficient Frontier with OSQP:\\
a Short Tutorial on the OR-Library Portfolio Benchmarks}
\author{Heuristic Algorithms in Transport and Finance --- Master's course assignment \\
    Javier Salvatierra \\
    Jesús Inurria
}
\date{September 2026}

\begin{document}
\maketitle
\vspace{-0.8cm}

\begin{abstract}
\noindent\small
The basic Markowitz portfolio selection problem is a convex quadratic program (QP): minimise the
variance of the portfolio return subject to a minimum expected return, a budget constraint and no
short selling. We formulate it, map it to the standard form of OSQP --- an open-source first-order
(ADMM) solver --- and trace the unconstrained efficient frontier of the five OR-Library instances
\texttt{port1}--\texttt{port5} by solving the 2000 QPs given by the return levels of
\texttt{portef1}--\texttt{portef5}. The frontiers match the reference ones to the rounding of the
benchmark files ($\sim10^{-5}\,\%$ in variance), at 0.4--137\,ms per QP. Near the maximum return the
problem becomes degenerate and 63 of the 10\,000 points are not solved to tolerance, which
illustrates the practical limits of a first-order method.
\end{abstract}

\section{Introduction}

An investor has a fixed budget and $N$ assets to spread it over. More expected return normally means
more uncertainty, so there is no single best portfolio but a trade-off. Markowitz~
made that trade-off computable by measuring risk as the \emph{variance} of the portfolio return, which
turns portfolio selection into an optimisation problem with a quadratic objective and linear
constraints. Beasley's tutorial~\cite{beasley2013} presents this basic model together with the
extensions that make it hard in practice (cardinality limits, transaction costs, index tracking); we
deal only with the basic version, without short selling.

The test data are the five portfolio instances of OR-Library~\cite{orlib}. It is built from weekly prices of real stocks between March 1992 and September
1997: Hang Seng, DAX 100, FTSE 100, S\&P 100 and Nikkei 225, with $N$ between 31 and 225 (see
Table~\ref{tab:results}). Each file \texttt{port$k$.txt} stores the mean return and the standard
deviation of every asset plus all pairwise correlations; the companion file \texttt{portef$k$.txt}
stores 2000 points (mean return, variance) of the corresponding \emph{unconstrained efficient
frontier} (UEF), which we use as the reference solution. The rest of the paper (i)~states the model,
(ii)~explains why we chose OSQP, and (iii)~reports times and accuracy on the five instances.

\section{The quadratic programming model}

\subsection{Parameters and decision variables}

For each asset $i=1,\dots,N$ the data give the expected (mean) return $\mu_i$ and the standard
deviation $s_i$ of its return, and for each pair the correlation $\rho_{ij}\in[-1,1]$. The decision
variable $w_i$ is the proportion of the budget invested in asset $i$, and $R_{\min}$ is the return
required from the portfolio. All statistics are weekly: $\mu_i=0.0109$ means $1.09\,\%$ per week.

\subsection{Formulation}

\begin{equation}
\min_{w_1,\dots,w_N}\ \sum_{i=1}^{N}\sum_{j=1}^{N} s_i\,s_j\,\rho_{ij}\,w_i w_j
\qquad\text{s.t.}\qquad
\sum_{i=1}^{N}\mu_i w_i \ = R_{\min},\qquad
\sum_{i=1}^{N} w_i = 1,\qquad
w_i \ge 0 .
\label{eq:model}
\end{equation}

Read literally: the objective is the variance of the portfolio return, obtained by adding, over every
pair of assets, the product of the two risks, their correlation and the two amounts invested. The
first constraint says that the portfolio must earn at least $R_{\min}$ on average. The second says
that the whole budget is invested. The third forbids short selling, which is caused by investing negative quantity any asset. The upper bound $w_i\le1$ need
not be written, since it follows from $w_i\ge0$ and $\sum_i w_i=1$.

Writing $\sigma_{ij}=s_i\rho_{ij}s_j$ for the covariance and $\Sigma=[\sigma_{ij}]$ for the
covariance matrix, the objective becomes the compact quadratic form $w^\top\Sigma w$. Beasley writes
the return constraint as an equality~\cite{beasley2013}; both versions give the same portfolio for
every $R_{\min}$ used here, because the constraint is active at the optimum (see below).

\subsection{Convexity}

A covariance matrix is always positive semidefinite, and in the five instances it is in fact positive
definite: the smallest eigenvalue ranges from $6.1\cdot10^{-6}$ (\texttt{port5}) to
$2.3\cdot10^{-4}$ (\texttt{port1}). Problem~\eqref{eq:model} is therefore a strictly convex QP: every
local minimum is global, the optimal portfolio is \emph{unique}, and the KKT conditions
\[
2\Sigma w-\lambda\mu-\nu\mathbf 1-s=0,\qquad \lambda,s\ge0,\qquad
\lambda\,(\mu^\top w-R_{\min})=0,\qquad s_iw_i=0
\]
are necessary and sufficient. This is what makes the model easy in the sense of Beasley: the objective
is non-linear, but efficient algorithms exist and the global optimum is found without difficulty. The
multiplier $\lambda$ is the shadow price of return, i.e.\ the slope $dV/dR$ of the frontier.

\subsection{The efficient frontier}

Solving~\eqref{eq:model} for one value of $R_{\min}$ gives one portfolio; sweeping $R_{\min}$ traces
the efficient frontier, the set of Pareto-optimal (non-dominated) portfolios. Its lower end is the
global minimum-variance portfolio (GVM), obtained by dropping the return constraint. The upper end is
$R_{\max}=\max_i\mu_i$, attained only by putting the whole budget in the single asset with the highest
mean return. The minimum variance $V^*(R)$ is convex,
non-decreasing and piecewise quadratic, each piece corresponding to a fixed set of assets held.
Individual assets, and any portfolio off the curve, are inefficient: they lie below and to the right of
it (red points in Figure~\ref{fig:frontiers}).

\section{Choice of solver}


Table~\ref{tab:solvers} lists the options suggested in the assignment. We chose
OSQP~\cite{osqpdocs}, an open-source solver written in C that applies ADMM~\cite{osqpdocs} to
problems in the standard form
\begin{equation}
\min_{w}\ \tfrac12 w^\top P w + q^\top w
\qquad\text{s.t.}\qquad l \le A w \le u ,
\qquad P\in\mathbb S^N_+ .
\label{eq:osqp}
\end{equation}
Model~\eqref{eq:model} maps onto~\eqref{eq:osqp} by taking $P=2\Sigma$ (the factor 2 cancels the
$\tfrac12$, so that the objective value returned is exactly the variance), $q=0$, and one row of $A$
per constraint: row~0 is $\mu^\top$ with $l_0=R_{\min}$ and $u_0=+\infty$, row~1 is $\mathbf 1^\top$
with $l_1=u_1=1$, and the remaining $N$ rows are the identity with $l=0$, $u=+\infty$.

\begin{table}[H]
\centering
\small
\caption{Candidate solvers for a convex QP.}
\label{tab:solvers}
\footnotesize
\begin{tabular}{@{}llll@{}}
\toprule
Solver & Algorithm for convex QP & Licence & Comment for this problem \\
\midrule
Gurobi / CPLEX & Barrier (interior point), simplex & Commercial (free academic) & Robust and accurate; also MIQP; need a licence \\
HiGHS & Active set & MIT & Accurate and lightweight; no MIQP \\
SCIP & Branch-and-bound for MINLP & Apache 2.0 & Oversized for a convex QP; useful with cardinality \\
\textbf{OSQP} & \textbf{ADMM / operator splitting} & \textbf{Apache 2.0} & \textbf{Chosen} \\
\bottomrule
\end{tabular}
\end{table}

The decisive argument is the \emph{structure of the task}: a frontier is not one QP but 2000 almost
identical QPs in which only the scalar $l_0=R_{\min}$ changes. Futhermore, it is used as a python libreary which makes the implementation easier.
\begin{itemize}
\item \textbf{One factorisation for the whole frontier.} The bounds $l,u$ do not enter the KKT matrix,
so \texttt{setup} factorises it once (5--12\,ms here) and \texttt{update(l=...)} changes the required
return without refactorising.
\item \textbf{Warm start.} Each solve restarts from the previous portfolio, which is very close to the
next one.
\item \textbf{Accuracy settings.} The user controls the tolerance, the iteration limit and
\emph{solution polishing}, a final linear solve on the detected active set that lifts an ADMM solution
to near machine precision. Iterations are cheap, there is no licence and installation is
\texttt{pip install osqp}.
\end{itemize}
The price to pay is that ADMM is a first-order method: it needs tuning (with the default tolerance
$10^{-3}$ and no polishing we measured variance errors of up to $3.1\,\%$, $11.2\,\%$ and $6.7\,\%$ on
\texttt{port1}--\texttt{port3}) and it converges slowly on degenerate problems (Section~\ref{sec:res}).

\section{Implementation}

The code is a Python notebook (\texttt{markowitz\_osqp.ipynb}, delivered separately) using Python
3.14, NumPy, SciPy, pandas, matplotlib and OSQP 1.1.3. It reads \texttt{port$k$.txt}, builds
$\Sigma=\rho\odot ss^\top$, sets the problem up once per instance and then loops over the return
levels:
\begin{center}
\texttt{for R in R\_levels: \quad l[0] = R; \quad solver.update(l=l); \quad res = solver.solve()}
\end{center}
The settings are \texttt{eps\_abs = eps\_rel = 1e-6}, \texttt{polishing = True} and
\texttt{max\_iter = 20000} (OSQP's default of 4000 is not enough near $R_{\max}$); the rest are left at
their defaults. As grid we use the \textbf{2000 return levels of \texttt{portef$k$}} themselves,
which allows a point-by-point comparison without interpolation. We record the wall-clock time per
instance (setup plus 2000 solves), the mean number of ADMM iterations, the points whose status is not
\texttt{solved}, and the relative deviation
$100\,(V_{\text{OSQP}}-V_{\text{portef}})/V_{\text{portef}}$ over the solved points. All runs are
single-threaded on an AMD Ryzen 5 5500U laptop under Windows 11. As a sanity check, \texttt{port1}
with $R_{\min}=0.008$ is solved in 125 iterations and 1.0\,ms; the portfolio holds 4 of the 31 assets
(40.1, 37.5, 16.7 and 5.7\,\%), returns exactly $0.008000$ and has variance $1.545024\cdot10^{-3}$.

\section{Computational results}
\label{sec:res}


\begin{figure}[!htb]
\centering
\includegraphics[width=0.83\textwidth]{frontiers.png}
\caption{Panels 1--5: unconstrained efficient frontier computed with OSQP (blue line), 26 reference
points of \texttt{portef} (black circles) and the individual assets (red dots), which are inefficient
portfolios. Panel 6: relative deviation in variance along the frontier, with position measured as
$(R-R_{\mathrm{GMV}})/(R_{\max}-R_{\mathrm{GMV}})$.}
\label{fig:frontiers}
\end{figure}

Table~\ref{tab:results} collects the results and Figure~\ref{fig:frontiers} shows the five frontiers.

\begin{table}[H]
\centering
\small
\caption{Tracing the five frontiers: 2000 QPs per instance, OSQP 1.1.3, single thread.}
\label{tab:results}
\footnotesize
\begin{tabular}{@{}llrrrrrrr@{}}
\toprule
Instance & Index & $N$ & Time (s) & ms/QP & Mean iter. & Not solved & Max $|$dev$|$ (\%) & Mean $|$dev$|$ (\%) \\
\midrule
\texttt{port1} & Hang Seng  & 31  & 0.87   & 0.44   & 76  & 0  & $7.8\cdot10^{-6}$ & $2.9\cdot10^{-6}$ \\
\texttt{port2} & DAX 100    & 85  & 25.93  & 12.96  & 797 & 0  & $3.6\cdot10^{-5}$ & $1.0\cdot10^{-5}$ \\
\texttt{port3} & FTSE 100   & 89  & 10.08  & 5.04   & 255 & 0  & $3.8\cdot10^{-2}$ & $3.6\cdot10^{-5}$ \\
\texttt{port4} & S\&P 100   & 98  & 33.90  & 16.95  & 803 & 3  & $1.0\cdot10^{-2}$ & $1.7\cdot10^{-5}$ \\
\texttt{port5} & Nikkei 225 & 225 & 274.96 & 137.48 & 974 & 60 & $1.1\cdot10^{-1}$ & $1.2\cdot10^{-4}$ \\
\bottomrule
\end{tabular}
\end{table}


\textbf{Accuracy.} Over almost the whole curve the agreement with the benchmark is of order
$10^{-5}\,\%$, which is simply the rounding of \texttt{portef} (10 decimals for variances of order
$10^{-4}$): It reproduces the reference frontiers exactly to the precision at which they are published.
The isolated spikes of up to $0.11\,\%$ in panel 6 are a handful of points where polishing failed; the
solution then stays at ADMM accuracy, violates the return constraint by $\sim10^{-8}$ and therefore
reports a slightly \emph{lower} variance.

\textbf{Computational time.} Tracing a whole frontier costs from 0.87\,s (\texttt{port1}) to about
4.5\,min (\texttt{port5}), roughly 5.8\,min for the 10\,000 QPs, i.e.\ 0.44--137\,ms per QP. Time is not
a function of $N$ alone --- \texttt{port3} ($N=89$) averages 255 iterations while \texttt{port2}
($N=85$) needs 797 --- but of conditioning and of the geometry of the upper end. Profiling shows that
$\sim96\,\%$ of the time is spent inside OSQP's C iterations and that the last $5\,\%$ of the return
levels consume $73\,\%$ of the time of \texttt{port5}; timings fluctuate between runs (183--275\,s).

\textbf{The failures at the top end.} 63 of the 10\,000 points are not solved to tolerance: 3 in
\texttt{port4} ($R\in[0.009090,0.009097]$, $R_{\max}=0.009195$) and 60 in \texttt{port5}
($R\in[0.003856,0.003971]$, 44 of them hitting the iteration limit). All lie next to $R_{\max}$, and the
reason is geometric rather than numerical: reaching a return close to $R_{\max}$ forces almost all the
budget into the single best asset, so the feasible set shrinks to a very thin corner of the simplex,
and at $R_{\max}$ exactly it is a single point with no strictly feasible portfolio. 
The problem is degenerate: at $R_{\max}$ only one portfolio is feasible, and the price of return
$dV/dR$ becomes huge. OSQP builds that price incrementally, in steps of size $\rho$, so it needs tens
of thousands of iterations to get there. With the default \texttt{max\_iter = 4000} the damage is much
larger (121, 82 and 125 lost points in \texttt{port2}, \texttt{port4} and \texttt{port5}), and some
feasible problems are even reported as \emph{primal infeasible}.
\vspace{-2pt}
\section{Conclusions}

The basic Markowitz model is a strictly convex QP, and an open-source first-order solver reproduces
the OR-Library reference frontiers to their published precision, provided the tolerance is tightened
to $10^{-6}$ and polishing is enabled. Exploiting the parametric structure --- one factorisation plus
warm starts --- brings each of the 10\,000 QPs down to milliseconds. The weak spot is the
maximum-return end of the frontier, which is degenerate and where a first-order method is not the right
tool: solver choice depends on problem geometry, not only on size.

\paragraph{Code availability.} All computations are reproducible from the notebook
\texttt{markowitz\_osqp.ipynb}, delivered separately: open it with the project kernel and \emph{Run
All} ($\approx6$\,min).

\begin{thebibliography}{9}
\footnotesize\setlength{\itemsep}{0pt}\setlength{\parskip}{0pt}
\bibitem{beasley2013} J.~E. Beasley. Portfolio optimisation: models and solution approaches. \emph{INFORMS TutORials in OR}, 201--221, 2013.
\bibitem{orlib} J.~E. Beasley. OR-Library. \emph{J.\ Oper.\ Res.\ Soc.}, 41(11):1069--1072, 1990. Data: \url{people.brunel.ac.uk/~mastjjb/jeb/orlib/portinfo.html}
\bibitem{osqpdocs} OSQP Developer Team. \emph{OSQP documentation: Python interface}, 2025. \url{https://osqp.org/docs/interfaces/python.html}
\end{thebibliography}

\end{document}
